\chapter{In Which Mazarin Becomes Prodigal}

\lettrine{W}{hilst} Mazarin was endeavouring to recover from the  serious alarm he had just experienced, Athos and Raoul were exchanging a few  words in a corner of the apartment. <Well, here you are at Paris, then,  Raoul?> said the comte.

\zz
<Yes, monsieur, since the return of M. le Prince.>

<I cannot converse freely with you here, because we are observed; but I  shall return home presently, and shall expect you as soon as your duty  permits.>

Raoul bowed, and, at that moment, M. le Prince came up to them. The prince had  that clear and keen look which distinguishes birds of prey of the noble  species; his physiognomy itself presented several distinct traits of this  resemblance. It is known that in the Prince de Condé, the aquiline nose rose  out sharply and incisively from a brow slightly retreating, rather low than  high, and according to the railers of the court,—a pitiless race without  mercy even for genius,—constituted rather an eagle's beak than a  human nose, in the heir of the illustrious princes of the house of Condé. This  penetrating look, this imperious expression of the whole countenance, generally  disturbed those to whom the prince spoke, more than either majesty or regular  beauty could have done in the conqueror of Rocroi. Besides this, the fire  mounted so suddenly to his projecting eyes, that with the prince every sort of  animation resembled passion. Now, on account of his rank, everybody at the  court respected M. le Prince, and many even, seeing only the man, carried their  respect as far as terror.

Louis de Condé then advanced towards the Comte de la Fère and Raoul, with the  marked intention of being saluted by the one, and of speaking with the other.  No man bowed with more reserved grace than the Comte de la Fère. He disdained  to put into a salutation all the shades which a courtier ordinarily borrows  from the same colour—the desire to please. Athos knew his own personal  value, and bowed to the prince like a man, correcting by something sympathetic  and undefinable that which might have appeared offensive to the pride of the  highest rank in the inflexibility of his attitude. The prince was about to  speak to Raoul. Athos forestalled him. <If M. le Vicomte de  Bragelonne,> said he, <were not one of the humble servants of your  royal highness, I would beg him to pronounce my name before you—mon  prince.>

<I have the honour to address Monsieur le Comte de la Fère,> said  Condé, instantly.

<My protector,> added Raoul, blushing.

<One of the most honourable men in the kingdom,> continued the  prince; <one of the first gentlemen of France, and of whom I have heard  so much that I have frequently desired to number him among my friends.>

<An honour of which I should be unworthy,> replied Athos, <but  for the respect and admiration I entertain for your royal highness.>

<Monsieur de Bragelonne,> said the prince, <is a good  officer, and it is plainly seen that he has been to a good school. Ah, monsieur  le comte, in your time, generals had soldiers!>

<That is true, my lord, but nowadays soldiers have generals.>

This compliment, which savoured so little of flattery, gave a thrill of joy to  the man whom already Europe considered a hero; and who might be thought to be  satiated with praise.

<I regret very much,> continued the prince, <that you should  have retired from the service, monsieur le comte; for it is more than probable  that the king will soon have a war with Holland or England, and opportunities  for distinguishing himself would not be wanting for a man who, like you, knows  Great Britain as well as you do France.>

<I believe I may say, monseigneur, that I have acted wisely in retiring  from the service,> said Athos, smiling. <France and Great Britain  will henceforward live like two sisters, if I can trust my  presentiments.>

<Your presentiments?>

<Stop, monseigneur, listen to what is being said yonder, at the table of  my lord the cardinal.>

<Where they are playing?>

<Yes, my lord.>

The cardinal had just raised himself on one elbow, and made a sign to the  king's brother, who went to him.

<My lord,> said the cardinal, <pick up, if you please, all  those gold crowns.> And he pointed to the enormous pile of yellow and  glittering pieces which the Comte de Guiche had raised by degrees before him by  a surprising run of luck at play.

<For me?> cried the Duc d'Anjou.

<Those fifty thousand crowns; yes, monseigneur, they are yours.>

<Do you give them to me?>

<I have been playing on your account, monseigneur,> replied the  cardinal, getting weaker and weaker, as if this effort of giving money had  exhausted all his physical and moral faculties.

<Oh, good heavens!> exclaimed Philip, wild with joy, <what a  fortunate day!> And he himself, making a rake of his fingers, drew a part  of the sum into his pockets, which he filled, and still full a third remained  on the table.

<Chevalier,> said Philip to his favourite, the Chevalier de  Lorraine, <come hither, chevalier.> The favourite quickly obeyed.  <Pocket the rest,> said the young prince.

This singular scene was considered by the persons present only as a touching  kind of family fete. The cardinal assumed the airs of a father with the sons of  France, and the two princes had grown up under his wing. No one then imputed to  pride, or even impertinence, as would be done nowadays, this liberality on the  part of the first minister. The courtiers were satisfied with envying the  prince.—The king turned away his head.

<I never had so much money before,> said the young prince,  joyously, as he crossed the chamber with his favourite to go to his carriage.  <No, never! What a weight these crowns are!>

<But why has monsieur le cardinal given away all this money at  once?> asked M. le Prince of the Comte de la Fère. <He must be very  ill, the dear cardinal!>

<Yes, my lord, very ill, without doubt; he looks very ill, as your royal  highness may perceive.>

<But surely he will die of it. A hundred and fifty thousand livres! Oh,  it is incredible! But, comte, tell me a reason for it?>

<Patience, monseigneur, I beg of you. Here comes M. le Duc d'Anjou,  talking with the Chevalier de Lorraine; I should not be surprised if they  spared us the trouble of being indiscreet. Listen to them.>

In fact the chevalier said to the prince in a low voice, <My lord, it is  not natural for M. Mazarin to give you so much money. Take care! you will let  some of the pieces fall, my lord. What design has the cardinal upon you to make  him so generous?>

<As I said,> whispered Athos in the prince's ear;  <that, perhaps, is the best reply to your question.>

<Tell me, my lord,> repeated the chevalier impatiently, as he was  calculating, by weighing them in his pocket, the quota of the sum which had  fallen to his share by rebound.

<My dear chevalier, a wedding present.>

<How a wedding present?>

<Eh! yes, I am going to be married,> replied the Duc d'Anjou,  without perceiving, at the moment, he was passing the prince and Athos, who  both bowed respectfully.

The chevalier darted at the young duke a glance so strange, and so malicious,  that the Comte de la Fère quite started on beholding it.

<You! you to be married!> repeated he; <oh! that's  impossible. You would not commit such a folly!>

<Bah! I don't do it myself; I am made to do it,> replied the  Duc d'Anjou. <But come, quick! let us get rid of our money.>  Thereupon he disappeared with his companion, laughing and talking, whilst all  heads were bowed on his passage.

<Then,> whispered the prince to Athos, <that is the  secret.>

<It was not I who told you so, my lord.>

<He is to marry the sister of Charles II?>

<I believe so.>

The prince reflected for a moment, and his eye shot forth one of its not  infrequent flashes. <Humph!> said he slowly, as if speaking to  himself; <our swords are once more to be hung on the wall—for a  long time!> and he sighed.

All that sigh contained of ambition silently stifled, of extinguished illusions  and disappointed hopes, Athos alone divined, for he alone heard that sigh.  Immediately after, the prince took leave and the king left the apartment.  Athos, by a sign made to Bragelonne, renewed the desire he had expressed at the  beginning of the scene. By degrees the chamber was deserted, and Mazarin was  left alone, a prey to suffering which he could no longer dissemble.  <Bernouin! Bernouin!> cried he in a broken voice.

<What does monseigneur want?>

<Guenaud—let Guenaud be sent for,> said his eminence.  <I think I'm dying.>

Bernouin, in great terror, rushed into the cabinet to give the order, and the  piqueur, who hastened to fetch the physician, passed the king's carriage  in the Rue Saint Honore.  